\section{静态经营收益与最优控制}

本阶段严格按照 benchmark 主稿中的原始设定，推导三类企业的静态经营模块。目标是得到后续 Bellman 方程所使用的优化后流收益。这里不改写模型，不扩张参数依赖，也不把优化后的对象提前写进 primitive objective。

\subsection{Benchmark Fidelity Check}

我们先重述本阶段必须推导的 benchmark 对象。

\paragraph{类型 $A$ 的 benchmark 对象。}
主稿写为
\[
\pi_A(z;w,M)
=
\pi_A^G(z;w)+\pi_A^{O,\mathrm{inc}}(z;w,M)-f_A
+x_A\Omega_A(z)-\frac{\chi_A}{\eta}x_A^\eta,
\qquad \eta>1,
\]
其中
\[
\Omega_A(z)\equiv \lambda_A(z)R_A(z).
\]
这里，$\pi_A(z;w,M)$ 是在静态优化之前的 benchmark 当期收益；$x_A$ 是内生的静态控制；$\Omega_A(z)$ 是由原始 commercialization-success rate $\lambda_A(z)$ 和原始 gross value $R_A(z)$ 构成的派生一单位价值；后面要推导的 $\Pi_A(z;w,M)$ 是优化后的静态收益，因此它是派生出来的静态结构对象，而不是 primitive object。

\paragraph{类型 $B$ 的 benchmark 对象。}
主稿先定义
\[
H_B(z;p_m,M)=\lambda_B(z,p_m,M)\bigl(R_B(z)-p_m\bigr)-\tau_M(M),
\]
然后写
\[
\pi_B(z;p_m,M)
=
-f_B+x_BH_B(z;p_m,M)-\frac{\chi_B}{\eta}x_B^\eta.
\]
因此，类型 $B$ 的 primitive objective 用的是 $x_BH_B$，不是 $x_B[H_B]_+$。正部算子必须从优化问题中推导出来，而不是预先放进目标函数。
这里，$\lambda_B(z,p_m,M)$ 是原始 reduced-form commercialization-success rate，$R_B(z)$ 是成功商业化的原始 gross value，$\tau_M(M)$ 是政策相关的商业化摩擦。$H_B(z;p_m,M)$ 是派生的单 idea 净价值。$\pi_B(z;p_m,M)$ 是优化前的 benchmark 当期收益。$x_B^*(z;p_m,M)$ 与 $\Pi_B(z;p_m,M)$ 则是在本阶段内端生并派生出来的对象。

\paragraph{类型 $C$ 的 benchmark 对象。}
主稿写为
\[
\pi_C(z;p_m,w,M)
=
\pi_C^G(z;w)+\pi_C^{O,\mathrm{inc}}(z;w,M)-f_C
+\max_{m\ge 0}\{p_m m-c(m,z;w)\}.
\]
这里，$c(m,z;w)$ 是原始静态 manufacturing cost function，$m$ 是内生的合格受托生产供给。benchmark 收益 $\pi_C(z;p_m,w,M)$ 在写法上把 within-period maximization block 部分展开出来。后面要推导的 $\Pi_C(z;p_m,w,M)$ 是 fully optimized static payoff。

\paragraph{本阶段的范围。}
本阶段要推导的是
\[
x_A^*(z),\quad \Pi_A(z;w,M),\quad
x_B^*(z;p_m,M),\quad \Pi_B(z;p_m,M),\quad
m^*(z;p_m,w),\quad \Pi_C(z;p_m,w,M).
\]

\paragraph{本阶段使用的全局记号与定义域。}
Phase 1 已经把 capability state space 固定为 $Z\subset \mathbb R_+$ 的非空紧集。凡是涉及价格方向的连续性或有界性，我们都在固定的非空紧价格集上讨论：
\[
\mathcal P_m\subset\mathbb R_+,\qquad \mathcal W\subset\mathbb R_+.
\]
因此，类型 $A$ 的紧定义域是
\[
K_A:=Z.
\]
类型 $B$ 的紧定义域是
\[
K_B:=Z\times \mathcal P_m.
\]
类型 $C$ 的紧定义域是
\[
K_C:=Z\times \mathcal P_m\times \mathcal W.
\]
我们还固定正部算子
\[
[u]_+:=\max\{u,0\}.
\]
最后，针对 type-$C$ 问题，我们记
\[
c_m(m,z;w):=\frac{\partial c(m,z;w)}{\partial m},
\qquad
c_{mm}(m,z;w):=\frac{\partial^2 c(m,z;w)}{\partial m^2}
\]
时，指的是相应导数存在的情形。

\subsection{类型 $A$ 的静态问题}

\paragraph{经济角色。}
类型 $A$ 企业是 integrated firms。它们在 benchmark 中既进行生产，也选择 original-drug idea-generation intensity。变量 $x_A$ 是静态控制；$\Pi_A(z;w,M)$ 是优化后的静态收益，后面会进入 Bellman 方程。

\paragraph{步骤 1：精确写出静态问题。}
固定 $(z,w,M)$。benchmark 的静态问题是
\begin{equation}
\max_{x_A\ge 0}
\left\{
\pi_A^G(z;w)+\pi_A^{O,\mathrm{inc}}(z;w,M)-f_A
+x_A\Omega_A(z)-\frac{\chi_A}{\eta}x_A^\eta
\right\}.
\label{eq:p2A_problem}
\end{equation}
由于前面三项与 $x_A$ 无关，因此等价于
\begin{equation}
\max_{x_A\ge 0}
\left\{
+x_A\Omega_A(z)-\frac{\chi_A}{\eta}x_A^\eta
\right\}.
\label{eq:p2A_problem_reduced}
\end{equation}

\paragraph{步骤 2：选择变量、状态与给定对象。}
选择变量是
\[
x_A\in[0,\infty).
\]
状态变量是
\[
z\in Z.
\]
在这个静态问题里，企业把下列对象视为给定：
\begin{enumerate}
    \item 工资或综合要素价格 $w$；
    \item 政策 regime $M$；
    \item 派生的一单位价值 $\Omega_A(z)$；
    \item 曲率参数 $\eta$；
    \item 成本尺度参数 $\chi_A$；
    \item benchmark 现金流项 $\pi_A^G(z;w)$、$\pi_A^{O,\mathrm{inc}}(z;w,M)$ 与 $f_A$。
\end{enumerate}
在当前问题中，只有 $x_A$ 是被选择的。其余对象都被当作参数。

\paragraph{步骤 3：使用的假设。}
\textbf{Benchmark 假设：} $\eta>1$、$\chi_A>0$，以及 $\Omega_A(z)=\lambda_A(z)R_A(z)$。  
\textbf{补充正则性假设：}
\begin{enumerate}
    \item[(A1)] 对每个固定的 $z$，$\Omega_A(z)$ 有限。
    \item[(A2)] 为了让 benchmark 闭式写法不出现正部算子，假设所有 $z$ 上都有 $\Omega_A(z)\ge 0$。
    \item[(A3)] 为了后面的连续性与有界性说明，假设在紧集 $Z$ 上，$z\mapsto \pi_A^G(z;w)$、$z\mapsto \pi_A^{O,\mathrm{inc}}(z;w,M)$ 与 $z\mapsto \Omega_A(z)$ 都连续且有界。
\end{enumerate}

\paragraph{步骤 4：证明最优解存在。}
定义
\[
g_A(x_A;z):=x_A\Omega_A(z)-\frac{\chi_A}{\eta}x_A^\eta.
\]
我们现在逐步证明简化问题 \eqref{eq:p2A_problem_reduced} 至少有一个最大值点。

首先，$g_A(\cdot;z)$ 在 $[0,\infty)$ 上连续，因为它是两个关于 $x_A$ 的连续函数之和：线性项 $x_A\Omega_A(z)$ 与幂函数项 $-(\chi_A/\eta)x_A^\eta$。

其次，证明大的 $x_A$ 不可能最优。由于 $\eta>1$，
\[
g_A(x_A;z)
=
+x_A\left[\Omega_A(z)-\frac{\chi_A}{\eta}x_A^{\eta-1}\right]\to -\infty
\qquad\text{当 }x_A\to\infty.
\]
原因是当 $\eta-1>0$ 时，$x_A^{\eta-1}\to\infty$，因此方括号里的负项最终压倒有限常数 $\Omega_A(z)$。

于是存在某个 $R_A(z)>0$，使得
\[
g_A(x_A;z)<g_A(0;z)=0
\qquad\text{对所有 }x_A\ge R_A(z).
\]
所以任何最大值点都必须落在区间
\[
K_A(z):=[0,R_A(z)]
\]
之内。这个集合非空、闭且有界，因此在 $\mathbb R$ 中是紧集。

现在调用 Weierstrass Extreme Value Theorem：一个定义在非空紧集上的连续实值函数一定能取到最大值和最小值。我们已经验证了它的全部前提：
\begin{enumerate}
    \item $K_A(z)$ 非空；
    \item $K_A(z)$ 紧；
    \item $g_A(\cdot;z)$ 在 $K_A(z)$ 上连续。
\end{enumerate}
因此 $g_A(\cdot;z)$ 在 $K_A(z)$ 上取到最大值。又因为在 $[0,\infty)$ 上的任何最大值点都必须位于 $K_A(z)$ 内，所以原始静态问题存在最优解。

\paragraph{步骤 5：一阶导数。}
对简化目标函数求导：
\[
\frac{\partial}{\partial x_A}
\left(
 x_A\Omega_A(z)-\frac{\chi_A}{\eta}x_A^\eta
\right)
=
\Omega_A(z)-\chi_Ax_A^{\eta-1}.
\]
因此一阶导数为
\begin{equation}
g_A'(x_A;z)=\Omega_A(z)-\chi_Ax_A^{\eta-1}.
\label{eq:p2A_derivative}
\end{equation}

\paragraph{步骤 6：二阶导数、凹性与充分性。}
对任意内点 $x_A>0$，
\[
g_A''(x_A;z)=-\chi_A(\eta-1)x_A^{\eta-2}<0.
\]
因此内点目标函数严格凹。

更一般地，当 $\eta>1$ 时，映射 $x_A\mapsto x_A^\eta$ 在 $[0,\infty)$ 上严格凸。把严格凸函数乘以正数 $\chi_A/\eta$ 仍保持严格凸，再乘以 $-1$ 就变成严格凹。再加上线性项 $x_A\Omega_A(z)$ 不会改变严格凹性。因此 $g_A(\cdot;z)$ 在整个可行集 $[0,\infty)$ 上严格凹。

这意味着至多只有一个最大值点。所以一旦我们找到一个满足相应一阶逻辑的可行点，它就必须是唯一的全局最优解。

\paragraph{步骤 7：求解最优策略。}
现在区分不同情形。

\textit{情形 1：$\Omega_A(z)=0$。}
若 $\Omega_A(z)=0$，则对所有可行的 $x_A\ge 0$，
\[
g_A(x_A;z)=0-\frac{\chi_A}{\eta}x_A^\eta\le 0.
\]
而在 $x_A=0$ 处，
\[
g_A(0;z)=0.
\]
因此对任意 $x_A>0$，
\[
g_A(x_A;z)=-\frac{\chi_A}{\eta}x_A^\eta<0=g_A(0;z).
\]
所以唯一最优解是
\[
x_A^*(z)=0.
\]

\textit{情形 2：$\Omega_A(z)>0$。}
若 $\Omega_A(z)>0$，则边界点 $x_A=0$ 处的一阶导数为
\[
g_A'(0;z)=\Omega_A(z)>0.
\]
所以边界点不可能最优：稍微向右移动会先提高目标函数。

另一方面，因为 $x_A^{\eta-1}\to\infty$ 当 $x_A\to\infty$，
\[
g_A'(x_A;z)=\Omega_A(z)-\chi_Ax_A^{\eta-1}\to -\infty.
\]
而 $g_A'(\cdot;z)$ 在 $x_A$ 上连续，因为它由连续函数组成。因此由 Intermediate Value Theorem 可知，存在至少一个严格正的点 $x_A^*(z)>0$，满足
\[
\Omega_A(z)-\chi_A\bigl(x_A^*(z)\bigr)^{\eta-1}=0,
\]
也就是
\[
g_A'(x_A^*(z);z)=0.
\]
这里，Intermediate Value Theorem 的前提成立，因为：
\begin{enumerate}
    \item $g_A'(\cdot;z)$ 连续；
    \item $g_A'(0;z)>0$；
    \item 对足够大的 $x_A$，$g_A'(x_A;z)<0$。
\end{enumerate}
由于 $g_A$ 严格凹，这个导数零点唯一，因此它就是唯一的最大值点。

现在一步一步解一阶条件：
\[
\chi_A\bigl(x_A^*(z)\bigr)^{\eta-1}=\Omega_A(z),
\]
这是把第二项移到右边。接下来除以正数 $\chi_A$：
\[
\bigl(x_A^*(z)\bigr)^{\eta-1}=\frac{\Omega_A(z)}{\chi_A}.
\]
再两边同时取 $1/(\eta-1)$ 次幂：
\[
x_A^*(z)=\left(\frac{\Omega_A(z)}{\chi_A}\right)^{\frac{1}{\eta-1}}.
\]
这一步是合法的，因为 $\eta-1>0$，并且在当前情形下 $\Omega_A(z)/\chi_A>0$。

把两种情形合并起来，benchmark 的策略函数就是
\begin{equation}
x_A^*(z)=\left(\frac{\Omega_A(z)}{\chi_A}\right)^{\frac{1}{\eta-1}}.
\label{eq:p2A_policy}
\end{equation}
在 A2 下，右端对所有 $z$ 都有定义，包括零点情形。

\paragraph{步骤 8：回代到目标函数，得到优化后收益。}
把 $x_A=x_A^*(z)$ 代回 benchmark 收益：
\[
\Pi_A(z;w,M)
=
\pi_A^G(z;w)+\pi_A^{O,\mathrm{inc}}(z;w,M)-f_A
+x_A^*(z)\Omega_A(z)-\frac{\chi_A}{\eta}\bigl(x_A^*(z)\bigr)^\eta.
\]
由一阶条件可知
\[
\chi_A\bigl(x_A^*(z)\bigr)^{\eta-1}=\Omega_A(z).
\]
两边同乘 $x_A^*(z)$，得
\[
\chi_A\bigl(x_A^*(z)\bigr)^\eta=x_A^*(z)\Omega_A(z).
\]
因此
\[
\frac{\chi_A}{\eta}\bigl(x_A^*(z)\bigr)^\eta
=
\frac{1}{\eta}x_A^*(z)\Omega_A(z).
\]
于是
\begin{align*}
\Pi_A(z;w,M)
&=
\pi_A^G(z;w)+\pi_A^{O,\mathrm{inc}}(z;w,M)-f_A
+x_A^*(z)\Omega_A(z)-\frac{1}{\eta}x_A^*(z)\Omega_A(z)\\
&=
\pi_A^G(z;w)+\pi_A^{O,\mathrm{inc}}(z;w,M)-f_A
+\frac{\eta-1}{\eta}x_A^*(z)\Omega_A(z).
\end{align*}
再代入 \eqref{eq:p2A_policy}：
\[
x_A^*(z)\Omega_A(z)
=
\left(\frac{\Omega_A(z)}{\chi_A}\right)^{\frac{1}{\eta-1}}\Omega_A(z)
=
\chi_A^{-\frac{1}{\eta-1}}\Omega_A(z)^{\frac{\eta}{\eta-1}}.
\]
因此
\begin{equation}
\Pi_A(z;w,M)
=
\pi_A^G(z;w)+\pi_A^{O,\mathrm{inc}}(z;w,M)-f_A
+\frac{\eta-1}{\eta}\chi_A^{-\frac{1}{\eta-1}}\Omega_A(z)^{\frac{\eta}{\eta-1}}.
\label{eq:p2A_payoff}
\end{equation}

\paragraph{步骤 9：连续性与有界性。}
现在严格说明连续性。

定义标量映射
\[
T_A(u):=\left(\frac{u}{\chi_A}\right)^{\frac{1}{\eta-1}},
\qquad u\in[0,\infty).
\]
因为 $\chi_A>0$，映射 $u\mapsto u/\chi_A$ 在 $[0,\infty)$ 上连续；而幂函数 $u\mapsto u^{1/(\eta-1)}$ 也在 $[0,\infty)$ 上连续。所以 $T_A$ 在 $[0,\infty)$ 上连续。

在 A2 与 A3 下，$\Omega_A$ 的值域落在 $[0,\infty)$，并且 $z\mapsto \Omega_A(z)$ 在紧集 $Z$ 上连续。于是策略函数
\[
x_A^*(z)=T_A(\Omega_A(z))
\]
是连续函数的复合，因此连续。

因为 $Z$ 紧，Weierstrass theorem 表明 $Z$ 上的每个连续实值函数都能取到最大值和最小值，因此一定有界。把这一结论应用到 $x_A^*(z)$，就知道策略函数在 $Z$ 上有界。

最后，收益公式 \eqref{eq:p2A_payoff} 是若干有界连续函数的和与复合：
\begin{enumerate}
    \item $z\mapsto \pi_A^G(z;w)$ 由 A3 可知连续且有界；
    \item $z\mapsto \pi_A^{O,\mathrm{inc}}(z;w,M)$ 由 A3 可知连续且有界；
    \item 因为 $\Omega_A(z)\ge 0$ 且幂函数在 $[0,\infty)$ 上连续，所以 $z\mapsto \Omega_A(z)^{\eta/(\eta-1)}$ 连续。
\end{enumerate}
因此 $\Pi_A(z;w,M)$ 在 $K_A=Z$ 上连续且有界。

\subsubsection*{What did we just do, and why does it matter?}
我们精确求解了 integrated firm 的单期研发问题。经济含义很直接：类型 $A$ 会提高研发强度，直到再多投入一单位 idea generation 的边际收益等于边际凸成本。由于问题严格凹，这个规则给出唯一最优解。得到的 $\Pi_A(z;w,M)$ 就是后续动态模型在比较企业是否以类型 $A$ 运营时使用的当期收益对象。

\paragraph{What has been established mathematically.}
我们已经证明，benchmark 中类型 $A$ 的静态问题存在唯一解，推导出 benchmark 策略函数 $x_A^*(z)$，并得到优化后流收益 $\Pi_A(z;w,M)$。

\paragraph{What assumptions were used.}
我们用到了 $\eta>1$、$\chi_A>0$、benchmark 中 $\Omega_A(z)=\lambda_A(z)R_A(z)$ 的定义，以及补充假设：$\Omega_A(z)$ 有限、弱非负，并在紧集 $Z$ 上连续且有界。

\paragraph{What remains to be derived next.}
下一步是推导 benchmark 中类型 $B$ 的静态问题，尤其要说明 $[H_B]_+$ 如何从优化问题中推导出来，而不是 primitive object。

\subsection{类型 $B$ 的静态问题}

\paragraph{经济角色。}
类型 $B$ 企业是 research-specialized firms。它们在 benchmark 中不进行内部生产。其静态控制变量是 idea-generation intensity $x_B$。对象 $H_B(z;p_m,M)$ 总结了一单位 original-drug idea 的预期净商业化价值。优化后收益 $\Pi_B(z;p_m,M)$ 是后续进入 Bellman 方程的静态对象。

\paragraph{步骤 1：精确写出静态问题。}
固定 $(z,p_m,M)$。benchmark 先定义
\begin{equation}
H_B(z;p_m,M)=\lambda_B(z,p_m,M)\bigl(R_B(z)-p_m\bigr)-\tau_M(M).
\label{eq:p2B_HB}
\end{equation}
然后精确写出当期收益为
\begin{equation}
\pi_B(z;p_m,M)
=
-f_B+x_BH_B(z;p_m,M)-\frac{\chi_B}{\eta}x_B^\eta.
\label{eq:p2B_payoff_primitive}
\end{equation}
因此，benchmark 的静态优化问题是
\begin{equation}
\max_{x_B\ge 0}
\left\{
-f_B+x_BH_B(z;p_m,M)-\frac{\chi_B}{\eta}x_B^\eta
\right\}.
\label{eq:p2B_problem}
\end{equation}
等价地，因为 $-f_B$ 与 $x_B$ 无关，可以写成
\begin{equation}
\max_{x_B\ge 0}
\left\{
+x_BH_B(z;p_m,M)-\frac{\chi_B}{\eta}x_B^\eta
\right\}.
\label{eq:p2B_problem_reduced}
\end{equation}

\paragraph{步骤 2：选择变量、状态与给定对象。}
选择变量是
\[
x_B\in[0,\infty).
\]
状态变量是
\[
z\in Z.
\]
在这个静态问题里，企业把下列对象视为给定：
\begin{enumerate}
    \item 合格受托生产价格 $p_m$；
    \item 政策 regime $M$；
    \item 派生的一单位净价值 $H_B(z;p_m,M)$；
    \item 曲率参数 $\eta$；
    \item 成本尺度参数 $\chi_B$；
    \item 固定经营成本 $f_B$。
\end{enumerate}
当前问题里只有 $x_B$ 是被选择的。$H_B(z;p_m,M)$ 不是 choice variable，而是从 benchmark 环境中继承下来的 reduced-form 价值对象。

\paragraph{步骤 3：使用的假设。}
\textbf{Benchmark 假设：} $\eta>1$、$\chi_B>0$，以及 \eqref{eq:p2B_HB} 中对 $H_B$ 的定义。  
\textbf{补充正则性假设：}
\begin{enumerate}
    \item[(B1)] 对每个固定的 $(z,p_m,M)$，$H_B(z;p_m,M)$ 有限。
    \item[(B2)] 为了后面的连续性与有界性说明，假设在相关紧域上 $(z,p_m)\mapsto H_B(z;p_m,M)$ 连续且有界。
\end{enumerate}

\paragraph{步骤 4：证明最优解存在。}
定义
\[
g_B(x_B;z,p_m,M):=x_BH_B(z;p_m,M)-\frac{\chi_B}{\eta}x_B^\eta.
\]
我们再次严格验证存在性。

首先，$g_B(\cdot;z,p_m,M)$ 在 $[0,\infty)$ 上连续，因为它是线性项 $x_BH_B(z;p_m,M)$ 与连续幂项 $-(\chi_B/\eta)x_B^\eta$ 的和。

其次，证明足够大的 $x_B$ 不可能最优。由于 $\eta>1$ 且 $H_B(z;p_m,M)$ 有限，
\[
g_B(x_B;z,p_m,M)
=
 x_B\left[H_B(z;p_m,M)-\frac{\chi_B}{\eta}x_B^{\eta-1}\right]\to -\infty
\qquad\text{当 }x_B\to\infty.
\]
所以存在某个 $R_B(z,p_m,M)>0$，使得
\[
g_B(x_B;z,p_m,M)<g_B(0;z,p_m,M)=0
\qquad\text{对所有 }x_B\ge R_B(z,p_m,M).
\]
因此任何最大值点都必须位于
\[
K_B(z,p_m,M):=[0,R_B(z,p_m,M)]
\]
之内。该集合非空、闭且有界，所以在 $\mathbb R$ 中紧。

我们再次使用 Weierstrass Extreme Value Theorem。其条件满足，因为：
\begin{enumerate}
    \item $K_B(z,p_m,M)$ 非空；
    \item $K_B(z,p_m,M)$ 紧；
    \item $g_B(\cdot;z,p_m,M)$ 在该集上连续。
\end{enumerate}
因此 $g_B(\cdot;z,p_m,M)$ 在 $K_B(z,p_m,M)$ 上取到最大值，原始静态问题存在最优解。

\paragraph{步骤 5：一阶导数。}
对简化后的目标函数求导：
\[
\frac{\partial}{\partial x_B}
\left(
 x_BH_B(z;p_m,M)-\frac{\chi_B}{\eta}x_B^\eta
\right)
=
H_B(z;p_m,M)-\chi_Bx_B^{\eta-1}.
\]
因此
\begin{equation}
g_B'(x_B;z,p_m,M)=H_B(z;p_m,M)-\chi_Bx_B^{\eta-1}.
\label{eq:p2B_derivative}
\end{equation}

\paragraph{步骤 6：二阶导数与充分性。}
对任意内点 $x_B>0$，
\[
g_B''(x_B;z,p_m,M)=-\chi_B(\eta-1)x_B^{\eta-2}<0.
\]
因此目标函数严格凹。

更一般地，$x_B\mapsto x_B^\eta$ 在 $\eta>1$ 时严格凸。把它乘以正数 $\chi_B/\eta$ 仍严格凸，再乘以 $-1$ 就成了严格凹。再加上线性项 $x_BH_B(z;p_m,M)$ 不会改变严格凹性。因此 $g_B(\cdot;z,p_m,M)$ 在整个可行集 $[0,\infty)$ 上严格凹。

严格凹意味着最大值点至多一个。因此，一旦找到一个可行的最大值点，它就必然是唯一的全局最优解。

\paragraph{步骤 7：推导策略函数，并解释 $[H_B]_+$ 的作用。}
这是 benchmark 的关键一步。primitive objective 里含有的是 $H_B$，不是 $[H_B]_+$。正部算子只会在解出最优解以后出现。

\textit{情形 1：$H_B(z;p_m,M)\le 0$。}
取任意可行的 $x_B\ge 0$。因为 $H_B(z;p_m,M)\le 0$，所以
\[
x_BH_B(z;p_m,M)\le 0.
\]
同时，
\[
-\frac{\chi_B}{\eta}x_B^\eta\le 0,
\]
且当 $x_B>0$ 时严格小于零。

因此，对所有可行 $x_B\ge 0$，
\[
g_B(x_B;z,p_m,M)=x_BH_B(z;p_m,M)-\frac{\chi_B}{\eta}x_B^\eta\le 0.
\]
在边界点 $x_B=0$，
\[
g_B(0;z,p_m,M)=0.
\]
所以任何严格正的选择都给出更低的目标值，而边界值最好。因此
\[
x_B^*(z;p_m,M)=0
\]
是唯一最优解。

\textit{情形 2：$H_B(z;p_m,M)>0$。}
这时边界点处的一阶导数是
\[
g_B'(0;z,p_m,M)=H_B(z;p_m,M)>0.
\]
所以 $x_B=0$ 不可能最优。

再看
\[
g_B'(x_B;z,p_m,M)=H_B(z;p_m,M)-\chi_Bx_B^{\eta-1}\to -\infty
\qquad\text{当 }x_B\to\infty,
\]
因为 $\eta-1>0$ 意味着 $x_B^{\eta-1}\to\infty$。

导数 $g_B'(\cdot;z,p_m,M)$ 在 $x_B$ 上连续。因此由 Intermediate Value Theorem 可知，存在至少一个严格正的点 $x_B^*(z;p_m,M)>0$ 使得
\[
g_B'(x_B^*(z;p_m,M);z,p_m,M)=0.
\]
该定理的前提满足，因为：
\begin{enumerate}
    \item $g_B'(\cdot;z,p_m,M)$ 连续；
    \item $g_B'(0;z,p_m,M)>0$；
    \item 对足够大的 $x_B$，$g_B'(x_B;z,p_m,M)<0$。
\end{enumerate}
又因为目标函数严格凹，这个导数零点唯一，因此它就是唯一最优解。

现在一步一步解一阶条件：
\[
H_B(z;p_m,M)-\chi_B\bigl(x_B^*(z;p_m,M)\bigr)^{\eta-1}=0.
\]
先把第二项移到右边：
\[
\chi_B\bigl(x_B^*(z;p_m,M)\bigr)^{\eta-1}=H_B(z;p_m,M).
\]
再除以正数 $\chi_B$：
\[
\bigl(x_B^*(z;p_m,M)\bigr)^{\eta-1}=\frac{H_B(z;p_m,M)}{\chi_B}.
\]
因为此时 $H_B(z;p_m,M)>0$，右端严格正，所以可以再取 $1/(\eta-1)$ 次幂，得到
\[
x_B^*(z;p_m,M)=\left(\frac{H_B(z;p_m,M)}{\chi_B}\right)^{\frac{1}{\eta-1}}.
\]

\textit{紧凑写法。}
现在先定义正部算子：
\[
[u]_+:=\max\{u,0\}.
\]
把两种情形合并，可写成
\begin{equation}
x_B^*(z;p_m,M)=\left(\frac{[H_B(z;p_m,M)]_+}{\chi_B}\right)^{\frac{1}{\eta-1}},
\label{eq:p2B_policy}
\end{equation}
并且可以直接核对这个紧凑式：
\begin{enumerate}
    \item 若 $H_B(z;p_m,M)\le 0$，则 $[H_B(z;p_m,M)]_+=0$，于是 \eqref{eq:p2B_policy} 给出 $x_B^*=0$；
    \item 若 $H_B(z;p_m,M)>0$，则 $[H_B(z;p_m,M)]_+=H_B(z;p_m,M)$，于是 \eqref{eq:p2B_policy} 给出前面推导出的内点闭式解。
\end{enumerate}
所以 $[H_B]_+$ 不是 primitive object，而是对分段最优解的紧凑记号。

\paragraph{步骤 8：回代到收益函数。}
如果 $H_B(z;p_m,M)\le 0$，则 $x_B^*=0$，因此
\[
\Pi_B(z;p_m,M)=-f_B.
\]
如果 $H_B(z;p_m,M)>0$，则
\[
\Pi_B(z;p_m,M)
=
-f_B+x_B^*H_B(z;p_m,M)-\frac{\chi_B}{\eta}(x_B^*)^\eta.
\]
由一阶条件可知
\[
\chi_B(x_B^*)^{\eta-1}=H_B.
\]
两边同乘 $x_B^*$，得到
\[
\chi_B(x_B^*)^\eta=x_B^*H_B.
\]
因此
\[
\frac{\chi_B}{\eta}(x_B^*)^\eta=\frac{1}{\eta}x_B^*H_B.
\]
于是
\begin{align*}
\Pi_B(z;p_m,M)
&=
-f_B+x_B^*H_B-\frac{1}{\eta}x_B^*H_B\\
&=
-f_B+\frac{\eta-1}{\eta}x_B^*H_B.
\end{align*}
再代入闭式策略函数：
\[
x_B^*H_B
=
\left(\frac{H_B}{\chi_B}\right)^{\frac{1}{\eta-1}}H_B
=
\chi_B^{-\frac{1}{\eta-1}}H_B^{\frac{\eta}{\eta-1}}.
\]
所以在 benchmark 的紧凑记号下，
\begin{equation}
\Pi_B(z;p_m,M)
=
-f_B+\frac{\eta-1}{\eta}\chi_B^{-\frac{1}{\eta-1}}[H_B(z;p_m,M)]_+^{\frac{\eta}{\eta-1}}.
\label{eq:p2B_optimized}
\end{equation}

\paragraph{步骤 9：连续性与有界性。}
定义标量映射
\[
T_B(u):=\left(\frac{[u]_+}{\chi_B}\right)^{\frac{1}{\eta-1}},
\qquad u\in\mathbb R.
\]
映射 $u\mapsto [u]_+$ 在全体 $\mathbb R$ 上连续；$u\mapsto u/\chi_B$ 连续，因为 $\chi_B>0$；幂函数 $u\mapsto u^{1/(\eta-1)}$ 在 $[0,\infty)$ 上连续。因此 $T_B$ 在 $\mathbb R$ 上连续。

在 B2 下，$(z,p_m)\mapsto H_B(z;p_m,M)$ 在紧域 $K_B=Z\times\mathcal P_m$ 上连续且有界。因此策略函数
\[
x_B^*(z;p_m,M)=T_B(H_B(z;p_m,M))
\]
是连续函数的复合，所以在 $K_B$ 上连续。

因为 $K_B$ 紧，Weierstrass theorem 表明 $K_B$ 上每个连续实值函数都能取到最大值和最小值，因此一定有界。把这一结论应用到 $x_B^*(z;p_m,M)$，就知道策略函数在 $K_B$ 上有界。

优化后收益 \eqref{eq:p2B_optimized} 也是 $K_B$ 上若干有界连续函数的和与复合：常数项 $-f_B$ 连续，而映射
\[
u\mapsto [u]_+^{\eta/(\eta-1)}
\]
在 $\mathbb R$ 上连续。因此 $\Pi_B(z;p_m,M)$ 在 $K_B$ 上连续且有界。

\subsubsection*{What did we just do, and why does it matter?}
我们在不改写 primitive objective 的前提下，解出了 benchmark 中 research-specialized firm 的静态问题。关键在于 $H_B(z;p_m,M)$ 的符号。如果商业化条件不好，使得 $H_B\le 0$，企业会把 idea-generation margin 完全关闭，选择零努力；如果商业化条件足够好，使得 $H_B>0$，企业就按照严格凹优化条件选择正的研发强度。这正是为什么 $[H_B]_+$ 应该出现在策略函数与优化后收益中，而不应该提前写进原始目标函数。

\paragraph{What has been established mathematically.}
我们已经证明，benchmark 中类型 $B$ 的静态问题存在唯一解，推导出两种情形下的解、紧凑的 $[H_B]_+$ 表达式，以及优化后收益 $\Pi_B(z;p_m,M)$。

\paragraph{What assumptions were used.}
我们使用了 $\eta>1$、$\chi_B>0$、benchmark 中 $H_B$ 的定义，以及补充假设：$H_B$ 有限，并且在紧域上连续且有界。

\paragraph{What remains to be derived next.}
下一步是推导 benchmark 中 manufacturing-specialized firm 的静态问题，使三类企业的优化后流收益全部就绪。

\subsection{类型 $C$ 的静态问题}

\paragraph{经济角色。}
类型 $C$ 企业是 manufacturing-specialized firms。它们在 benchmark 中不产生新的 original-drug ideas。其静态控制变量是合格受托生产供给 $m$。优化后收益 $\Pi_C(z;p_m,w,M)$ 将在后续 Bellman 方程中作为 manufacturing specialist 的当期收益。

\paragraph{步骤 1：精确写出静态问题。}
固定 $(z,p_m,w,M)$。benchmark 主稿写为
\begin{equation}
\pi_C(z;p_m,w,M)
=
\pi_C^G(z;w)+\pi_C^{O,\mathrm{inc}}(z;w,M)-f_C
+\max_{m\ge 0}\left\{p_m m-c(m,z;w)\right\}.
\label{eq:p2C_benchmark_payoff}
\end{equation}
因此，其静态控制问题就是
\begin{equation}
\max_{m\ge 0}\left\{p_m m-c(m,z;w)\right\}.
\label{eq:p2C_problem}
\end{equation}
一旦这个问题被解出，优化后的 benchmark 收益就是
\begin{equation}
\Pi_C(z;p_m,w,M)
=
\pi_C^G(z;w)+\pi_C^{O,\mathrm{inc}}(z;w,M)-f_C
+\max_{m\ge 0}\left\{p_m m-c(m,z;w)\right\}.
\label{eq:p2C_optimized_definition}
\end{equation}

\paragraph{步骤 2：选择变量、状态与给定对象。}
选择变量是
\[
m\in[0,\infty).
\]
状态变量是
\[
z\in Z.
\]
在这个问题里给定的是：
\begin{enumerate}
    \item 合格受托生产价格 $p_m$；
    \item 工资或综合要素价格 $w$；
    \item 政策 regime $M$；
    \item benchmark 当期现金流项 $\pi_C^G(z;w)$、$\pi_C^{O,\mathrm{inc}}(z;w,M)$ 与 $f_C$；
    \item 原始 manufacturing cost function $c(m,z;w)$。
\end{enumerate}
当前问题里只有 $m$ 是被选择的。

\paragraph{步骤 3：使用的假设。}
\textbf{Benchmark 假设：} benchmark 中类型 $C$ 的收益结构 \eqref{eq:p2C_benchmark_payoff}。  
\textbf{补充正则性假设：}
\begin{enumerate}
    \item[(C1)] 对每个固定的 $(z,w)$，函数 $m\mapsto c(m,z;w)$ 在 $[0,\infty)$ 上连续。
    \item[(C2)] 对每个固定的 $(z,w)$，函数 $m\mapsto c(m,z;w)$ 关于 $m$ 连续可导且严格凸。
    \item[(C3)] 对每个固定的 $(z,w)$，成本函数满足 coercive：
    \[
    \lim_{m\to\infty}\frac{c(m,z;w)}{m}=\infty.
    \]
    \item[(C4)] 为了连续性说明，存在一个连续下界函数 $\underline c:[0,\infty)\to\mathbb R$，满足
    \[
    c(m,z;w)\ge \underline c(m)
    \qquad\text{对所有 }(m,z,w)\in[0,\infty)\times Z\times\mathcal W,
    \]
    且
    \[
    \lim_{m\to\infty}\frac{\underline c(m)}{m}=\infty.
    \]
    \item[(C5)] 目标映射
    \[
    (m,z,p_m,w)\mapsto p_m m-c(m,z;w)
    \]
    在 $[0,\infty)\times K_C$ 上联合连续。
\end{enumerate}

假设 C3 足以处理固定参数下的存在性；假设 C4 则是更强的，用来在整个紧参数集上构造统一的紧选择集。后者是后面干净调用 maximum theorem 的关键。

\paragraph{步骤 4：证明最优解存在。}
定义
\[
g_C(m;z,p_m,w):=p_m m-c(m,z;w).
\]
我们先对固定的参数三元组 $(z,p_m,w)$ 证明存在性。

由 C1，$m\mapsto g_C(m;z,p_m,w)$ 在 $[0,\infty)$ 上连续。困难在于可行集不是紧的，所以和前两个问题一样，我们先证明当 $m$ 很大时目标函数一定很小。

由 C3，对每个固定的 $(z,p_m,w)$，存在 $R_C(z,p_m,w)>0$，使得当 $m\ge R_C(z,p_m,w)$ 时，
\[
\frac{c(m,z;w)}{m}>p_m+1.
\]
两边乘以正数 $m$，得
\[
c(m,z;w)>(p_m+1)m.
\]
于是对所有 $m\ge R_C(z,p_m,w)$，
\[
g_C(m;z,p_m,w)=p_m m-c(m,z;w)<p_m m-(p_m+1)m=-m.
\]
所以
\[
g_C(m;z,p_m,w)\to -\infty
\qquad\text{当 }m\to\infty.
\]
因此任何最优解都必须位于有界区间
\[
K_C(z,p_m,w):=[0,R_C(z,p_m,w)].
\]
这个集合非空、闭且有界，因此紧。

现在再调用一次 Weierstrass Extreme Value Theorem。它的前提满足，因为：
\begin{enumerate}
    \item $K_C(z,p_m,w)$ 非空；
    \item $K_C(z,p_m,w)$ 紧；
    \item $g_C(\cdot;z,p_m,w)$ 在 $K_C(z,p_m,w)$ 上连续。
\end{enumerate}
于是 $g_C(\cdot;z,p_m,w)$ 在 $K_C(z,p_m,w)$ 上取到最大值，所以原始静态问题存在最优解。

\paragraph{步骤 5：一阶导数。}
对于内点 $m>0$，对目标函数求导：
\[
\frac{\partial}{\partial m}\left(p_m m-c(m,z;w)\right)=p_m-c_m(m,z;w).
\]
因此一阶导数是
\begin{equation}
g_C'(m;z,p_m,w)=p_m-c_m(m,z;w).
\label{eq:p2C_derivative}
\end{equation}

\paragraph{步骤 6：二阶导数与充分性。}
再求一次导数：
\[
g_C''(m;z,p_m,w)=-c_{mm}(m,z;w).
\]
由 C2，$m\mapsto c(m,z;w)$ 严格凸。因此 $m\mapsto -c(m,z;w)$ 严格凹；再加上线性项 $m\mapsto p_m m$ 仍保持严格凹。所以 $g_C(\cdot;z,p_m,w)$ 在 $[0,\infty)$ 上严格凹。

这意味着最大值点至多一个。因此，一旦我们识别出一个可行最大值点，唯一性就立刻成立。

\paragraph{步骤 7：通过显式分情况推导策略函数。}
因为可行集是半轴 $[0,\infty)$，所以必须区分边界情形与内点情形。

\textit{情形 1：边界解。}
若
\[
p_m\le c_m(0,z;w),
\]
则由严格凸性可知，边际成本函数 $m\mapsto c_m(m,z;w)$ 弱递增。因此对所有 $m\ge 0$，
\[
c_m(m,z;w)\ge c_m(0,z;w)\ge p_m.
\]
移项可得
\[
g_C'(m;z,p_m,w)=p_m-c_m(m,z;w)\le 0
\qquad\text{对所有 }m\ge 0.
\]
所以目标函数在整个可行集上弱递减，其最大值只能出现在左端点：
\begin{equation}
m^*(z;p_m,w)=0.
\label{eq:p2C_boundary}
\end{equation}

\textit{情形 2：内点解。}
若
\[
p_m>c_m(0,z;w),
\]
则
\[
g_C'(0;z,p_m,w)=p_m-c_m(0,z;w)>0.
\]
因为 $g_C'$ 连续，所以存在某个 $\delta>0$，使得
\[
g_C'(m;z,p_m,w)>0
\qquad\text{对所有 }m\in[0,\delta].
\]
这意味着
\[
g_C(\delta;z,p_m,w)>g_C(0;z,p_m,w),
\]
所以边界点 $m=0$ 不可能最优。

我们已经证明最优解存在。既然最优解不是边界点，它就必须是某个内点 $m^*(z;p_m,w)>0$。

这时调用 Fermat's theorem for interior optima：如果一个函数在开区间上可导，并且在内点取得局部最大值，那么该点的一阶导数必须等于零。这里前提满足，因为：
\begin{enumerate}
    \item $g_C(\cdot;z,p_m,w)$ 在 $(0,\infty)$ 上可导，由 C2 保证；
    \item 当前情形下最大值点是内点。
\end{enumerate}
因此内点最优解必须满足
\begin{equation}
c_m(m^*(z;p_m,w),z;w)=p_m.
\label{eq:p2C_interior}
\end{equation}
最后，由于 $c(\cdot,z;w)$ 严格凸，边际成本函数 $c_m(\cdot,z;w)$ 严格递增，所以 \eqref{eq:p2C_interior} 至多只有一个解。结合前面已经证明的存在性，这个解就是唯一最优解。

\paragraph{步骤 8：如果可以闭式，就写出策略函数。}
在一般 benchmark 下，最优策略只能由 \eqref{eq:p2C_boundary} 与 \eqref{eq:p2C_interior} 隐式表征。只有当进一步指定 $c(m,z;w)$ 的函数形式时，才可能得到闭式解。由于 benchmark 主稿在这里没有施加特殊函数形式，所以精确的 benchmark 结果就是下面这个分情形表征：
\begin{equation}
m^*(z;p_m,w)=
\begin{cases}
0, & \text{若 } p_m\le c_m(0,z;w),\\[0.4em]
\text{满足 } c_m(m,z;w)=p_m \text{ 的唯一 } m>0, & \text{若 } p_m>c_m(0,z;w).
\end{cases}
\label{eq:p2C_policy}
\end{equation}

\paragraph{步骤 9：回代到收益函数。}
将最优解代回 \eqref{eq:p2C_benchmark_payoff}。为了书写清楚，先定义辅助价值对象
\[
\Phi_C(z;p_m,w):=\max_{m\ge 0}\left\{p_m m-c(m,z;w)\right\}.
\]
这只是辅助记号，不会替代 benchmark 对象 $\Pi_C$；它只是给优化后的 contract-manufacturing block 起一个名字，方便书写代数。

按照最大值点的定义，
\[
\Phi_C(z;p_m,w)=p_m m^*(z;p_m,w)-c(m^*(z;p_m,w),z;w).
\]
因此优化后的 benchmark 收益是
\begin{equation}
\Pi_C(z;p_m,w,M)
=
\pi_C^G(z;w)+\pi_C^{O,\mathrm{inc}}(z;w,M)-f_C
+\Phi_C(z;p_m,w).
\label{eq:p2C_payoff}
\end{equation}
把辅助记号展开以后，就是
\[
\Pi_C(z;p_m,w,M)
=
\pi_C^G(z;w)+\pi_C^{O,\mathrm{inc}}(z;w,M)-f_C
+p_m m^*(z;p_m,w)-c(m^*(z;p_m,w),z;w).
\]
在不额外指定特殊成本函数的前提下，这个表达已经不能再一般性地继续化简了。所以 \eqref{eq:p2C_payoff} 就是精确的 benchmark 优化后收益。

\paragraph{步骤 10：连续性与有界性。}
现在说明策略函数与优化后收益的正则性。

\textit{Step 10a：在紧参数域上构造统一的紧选择集。}
令
\[
\bar p_m:=\max\mathcal P_m.
\]
因为 $\mathcal P_m$ 非空且紧，所以这个最大值存在且有限。

再定义
\[
\bar c_0:=\max_{(z,w)\in Z\times \mathcal W} c(0,z;w).
\]
这个最大值也存在且有限，因为 $(z,w)\mapsto c(0,z;w)$ 在紧集 $Z\times \mathcal W$ 上连续。

由 C4，存在连续下界函数 $\underline c$，使得
\[
\frac{\underline c(m)}{m}\to\infty
\qquad\text{当 }m\to\infty.
\]
于是存在某个 $\bar R_C>0$，使得当 $m\ge \bar R_C$ 时，
\[
\frac{\underline c(m)}{m}>\bar p_m+1.
\]
如果需要，还可以把 $\bar R_C$ 再取大一些，使得
\[
\bar R_C>\bar c_0.
\]
把上式乘以 $m>0$，得到
\[
\underline c(m)>(\bar p_m+1)m.
\]
现在取任意参数三元组 $(z,p_m,w)\in K_C$ 和任意 $m\ge \bar R_C$。因为 $p_m\le \bar p_m$ 且 $c(m,z;w)\ge \underline c(m)$，所以
\begin{align*}
g_C(m;z,p_m,w)
&=p_m m-c(m,z;w)\\
&\le \bar p_m m-\underline c(m)\\
&<\bar p_m m-(\bar p_m+1)m\\
&=-m\\
&<-\bar c_0.
\end{align*}
在边界点 $m=0$，
\[
g_C(0;z,p_m,w)=-c(0,z;w)\ge -\bar c_0.
\]
所以任何满足 $m\ge \bar R_C$ 的可行选择，其值都严格低于 $m=0$ 的值。因此最优解不可能满足 $m\ge \bar R_C$。

于是我们可以把原始可行集 $[0,\infty)$ 替换成紧区间
\[
\Gamma_C(z,p_m,w):=[0,\bar R_C],
\]
而不改变最优解。这个替换对整个紧参数集 $K_C$ 都成立。

\textit{Step 10b：应用 Berge Maximum Theorem。}
现在调用标准 compact-valued 版本的 Berge Maximum Theorem。其内容是：如果目标函数连续，且 feasible correspondence 连续并具有非空紧值，那么 value function 连续，而 argmax correspondence 非空、紧值且 upper hemicontinuous。

它的假设在这里都满足，因为：
\begin{enumerate}
    \item 由 C5，目标映射
    \[
    (m,z,p_m,w)\mapsto p_m m-c(m,z;w)
    \]
    在 $[0,\infty)\times K_C$ 上连续，因此在紧集 $[0,\bar R_C]\times K_C$ 上也连续；
    \item feasible correspondence $\Gamma_C(z,p_m,w)=[0,\bar R_C]$ 不随 $(z,p_m,w)$ 变化，因此作为 correspondence 是连续的；
    \item 每个 $\Gamma_C$ 的值都非空且紧。
\end{enumerate}
因此 value function
\[
(z,p_m,w)\mapsto \Phi_C(z;p_m,w)
\]
在 $K_C$ 上连续，而 argmax correspondence 非空、紧值且 upper hemicontinuous。

\textit{Step 10c：利用唯一性得到连续策略函数。}
我们已经证明，对每个 $(z,p_m,w)\in K_C$，最优解唯一。于是 argmax correspondence 在每个参数点上都是单点值。Berge theorem 的标准唯一性推论告诉我们，这样的唯一 selector 是连续的。因此策略函数
\[
(z,p_m,w)\mapsto m^*(z;p_m,w)
\]
在 $K_C$ 上连续。

\textit{Step 10d：优化后收益有界。}
因为 $K_C$ 紧，而 $\Pi_C(z;p_m,w,M)$ 在 $K_C$ 上连续，所以由 Weierstrass theorem 可知，$\Pi_C$ 在 $K_C$ 上取到最大值和最小值。因此它在 $K_C$ 上有界。于是，在 benchmark 加上 C1--C5 之后，类型 $C$ 的优化后收益在相关紧状态—价格域上连续且有界。

\subsubsection*{What did we just do, and why does it matter?}
我们精确求解了 manufacturing specialist 的单期问题。经济含义是：类型 $C$ 会比较 contract-manufacturing price $p_m$ 与边际成本。如果价格在零产出点就低于边际成本，企业不会供给任何受托生产；如果价格足够高，企业就会扩大供给，直到价格等于边际成本。严格凸性保证这个条件刻画的是唯一最优解。得到的 $\Pi_C(z;p_m,w,M)$ 就是后续动态模型比较 manufacturing specialization 与其他组织形式时所用的当期收益对象。

\paragraph{What has been established mathematically.}
我们已经证明，benchmark 中类型 $C$ 的静态问题存在唯一解，得到 $m^*(z;p_m,w)$ 的精确分情形刻画，并获得优化后流收益 $\Pi_C(z;p_m,w,M)$。

\paragraph{What assumptions were used.}
我们使用了 benchmark 中类型 $C$ 的收益结构，以及补充假设：$c(m,z;w)$ 连续、关于 $m$ 连续可导、严格凸、对固定参数 coercive、在紧参数域上有统一下界，并且与状态和价格联合连续。

\paragraph{What remains to be derived next.}
至此，三类企业的 benchmark 静态经营模块都已完成。下一阶段是把 $\Pi_A(z;w,M)$、$\Pi_B(z;p_m,M)$ 与 $\Pi_C(z;p_m,w,M)$ 放进 Bellman 系统，并证明相应的动态性质。

\subsection{Phase 2 结果总结}

本阶段已经推导出进入动态系统的三个 benchmark 优化后流收益对象：
\[
\Pi_A(z;w,M),\qquad \Pi_B(z;p_m,M),\qquad \Pi_C(z;p_m,w,M).
\]
对类型 $A$，我们得到了唯一的 benchmark 研发强度策略函数及其优化后收益。对类型 $B$，我们从 primitive benchmark objective 出发推导了策略函数，并说明了为什么正部算子 $[H_B]_+$ 是由优化结果产生的紧凑记号，而不是 primitive ingredient。对类型 $C$，我们通过显式的边界 versus interior 分析推导出唯一的受托生产供给规则，并得到 benchmark 中精确的优化后收益。

这三个对象现在都可以直接代入 Bellman 方程，而不需要再做任何额外的静态改写。

\subsection{本阶段使用的假设}

本阶段使用了 benchmark 中三类企业的收益结构，以及补充正则性条件：类型 $A$ 与类型 $B$ 的创新成本凸性、类型 $C$ 的制造成本连续严格凸且 coercive、reduced-form 价值对象有限，以及相关 primitive 函数在紧域上的连续性与有界性。

\subsection{下一阶段的任务}

Phase 3 将证明 $\mathcal C_b(Z)^3$ 上的 Bellman 算子性质：
\begin{enumerate}
    \item 从有界连续函数三元组映射到有界连续函数三元组；
    \item 在 sup norm 下构成压缩映射；
    \item $(V_A,V_B,V_C)$ 的存在性与唯一性；
    \item 最优组织选择策略的存在，并严格沿用 benchmark 的 choice objects，而不是改写后的静态模板。
\end{enumerate}
